% =============================================================================
% Seção 1 — Introdução
% =============================================================================
\section{Introdução}
\label{sec:introducao}

As fundações compõem o subsistema estrutural responsável por transferir as ações da superestrutura ao maciço de solo de apoio, exercendo papel determinante na estabilidade global, na limitação de recalques diferenciais e na viabilidade econômica das edificações \citepabnt{abnt6122}. Em especial, as fundações superficiais do tipo sapata isolada apresentam ampla aplicação em edificações de pequeno e médio porte, em razão da simplicidade construtiva e do baixo custo associado \citepabnttwo{wang2008economic}{waheed2022practical}.

A despeito de sua frequência em obras correntes, o dimensionamento dessas fundações é, na prática profissional, frequentemente conduzido por procedimentos iterativos de tentativa e verificação, nos quais o engenheiro arbitra dimensões e verifica, sucessivamente, o atendimento a critérios de projeto associados à tensão admissível do solo, à punção e à compatibilidade geométrica entre o pilar e a sapata \citepabntthree{wang2008economic}{khajehzadeh2012gsa}{gandomi2018construction}. Esse fluxo de trabalho depende da experiência do projetista e pode produzir soluções viáveis, porém conservadoras, com consumo de concreto superior ao obtido por formulações de otimização econômica \citepabnttwo{waheed2022practical}{waheed2025optimization}.

A elevada pegada de carbono associada à produção do cimento e à execução das fundações reforça a relevância da racionalização do volume de concreto empregado. Ferramentas computacionais que automatizem a busca por soluções de menor consumo material podem, portanto, contribuir para a sustentabilidade das obras de engenharia civil, desde que preservados os critérios de segurança e desempenho \citepabnt{waheed2025optimization}. Nesse contexto, métodos de otimização baseados em metaheurísticas e em técnicas de aprendizado de máquina vêm sendo aplicados a problemas de projeto estrutural e geotécnico, dada a sua capacidade de explorar espaços de busca não convexos, de elevada dimensionalidade e com restrições heterogêneas \citepabntfour{gomes2018probabilistic}{morales2020balance}{abualigah2021arithmetic}{kashani2020optimum}.

Este estudo se insere nesse cenário e investiga a seguinte pergunta de pesquisa: \textit{é possível desenvolver uma abordagem computacional que realize o pré-dimensionamento geométrico e, em etapa posterior, o posicionamento otimizado de fundações superficiais, integrando verificações estruturais e geotécnicas em um único ambiente, de forma a reduzir o consumo de concreto e o esforço manual do projetista?}

Como recorte específico deste artigo, apresentam-se os resultados relativos à \textit{primeira} frente da pesquisa: o pré-dimensionamento geométrico otimizado das sapatas isoladas, considerando posições de pilares fornecidas pelo projeto estrutural e um conjunto limitado de verificações de tensão no solo, geometria e punção. A abordagem adotada é a otimização global assistida por modelo substituto, mais especificamente o método \textit{Efficient Global Optimization} (EGO) proposto por \citeonline{jones1998efficient}, no qual um Processo Gaussiano (GPR) atua como modelo substituto da função objetivo \citepabntthree{williams1995gaussian}{schulz2018tutorial}{shahriari2016review}, e um Algoritmo Genético (AG) é empregado como otimizador interno na maximização da função de aquisição \textit{Expected Improvement} (EI). A frente complementar -- a incorporação do problema de empacotamento ao posicionamento conjunto das fundações -- é objeto da próxima etapa do plano de trabalho e não é tratada de forma exaustiva neste manuscrito.

Como a função pseudo-objetivo desta etapa tem baixo custo computacional na implementação vetorizada atual, a adequação da arquitetura assistida por modelo substituto não é assumida: ela é avaliada empiricamente neste artigo sob dois regimes de orçamento — eficiência amostral com número igual de avaliações reais e tempo de parede com orçamento estendido para as buscas diretas — e complementada por uma auditoria de decomposição por sapata. Com isso, delimita-se o domínio em que o EGO é a escolha vantajosa, o domínio em que buscas diretas são suficientes e o limite metodológico imposto pela quase separabilidade dos casos atuais.

\subsection{Objetivos}
\label{sec:objetivos}

\paragraph{Objetivo geral.} Apresentar e discutir, do ponto de vista metodológico e experimental, uma abordagem computacional assistida por modelo substituto para o pré-dimensionamento geométrico de sapatas isoladas, integrando verificações parciais da mecânica das estruturas e da geotecnia, e implementada em uma plataforma computacional (\fundaai{}) construída em Python.

\paragraph{Objetivos específicos.}
\begin{enumerate}
    \item Revisar criticamente o estado da arte recente em otimização aplicada ao projeto de fundações superficiais, com ênfase em metaheurísticas e em métodos baseados em modelos substitutos.
    \item Formular o problema de pré-dimensionamento geométrico de sapatas isoladas como uma otimização mono-objetivo penalizada, combinando critérios inspirados na NBR~6118 \citepabnt{abnt6118} e na NBR~6122 \citepabnt{abnt6122}, correlações empíricas e hipóteses explícitas de pré-dimensionamento.
    \item Apresentar os fundamentos teóricos do método EGO e do GPR e discutir criticamente a escolha dessa arquitetura híbrida em uma formulação que, nesta etapa, possui avaliação direta de baixo custo, mas prepara o método para extensões com maior custo de avaliação e restrições mais complexas.
    \item Reportar, sob protocolo experimental fechado (sementes registradas, orçamento de avaliações equalizado e teste estatístico pareado), o efeito do fator de penalidade e da escolha do \textit{kernel} sobre a qualidade preditiva do modelo substituto, a comparação da arquitetura proposta com Algoritmo Genético, PSO, GWO, busca aleatória e uma variante com tratamento explícito de restrições (\textit{Constrained Bayesian Optimization}) em dois regimes de orçamento, e uma auditoria de quase separabilidade por decomposição de sapatas.
    \item Discutir as limitações da abordagem atual e situar, como trabalho subsequente, a incorporação explícita do problema de \textit{empacotamento} para o posicionamento conjunto e o dimensionamento das fundações.
\end{enumerate}

\subsection{Contribuição e estrutura do artigo}
\label{sec:contribuicao}

A principal contribuição deste artigo é apresentar, em formato consolidado, uma arquitetura EGO--GPR--AG aplicada à minimização do volume de concreto de sapatas isoladas e discutir as escolhas de modelagem que afetam a qualidade preditiva do modelo substituto -- em particular, o fator de penalidade aplicado às restrições e a função de covariância do GPR --, quantificando ainda, sob protocolo controlado, quando o tratamento explícito das restrições por otimização bayesiana com restrições supera a penalização direta. O artigo também inclui uma linha de base por decomposição para explicitar o quanto os casos atuais são quase separáveis. Adicionalmente, o trabalho documenta a implementação de uma plataforma computacional (\fundaai{}) que integra leitura de dados de projeto, parametrização do otimizador, execução do método e visualização dos arranjos resultantes, com todos os resultados regeneráveis por \textit{scripts} determinísticos a partir de artefatos persistidos.

O restante do artigo está organizado da seguinte forma. A Seção~\ref{sec:estado_da_arte} apresenta o estado da arte sobre otimização aplicada ao dimensionamento de fundações superficiais e o panorama metodológico das técnicas baseadas em modelos substitutos. A Seção~\ref{sec:fundamentacao} introduz a fundamentação teórica do EGO, do GPR e da função de aquisição EI. A Seção~\ref{sec:metodologia} formaliza o problema de otimização e descreve a arquitetura proposta. A Seção~\ref{sec:implementacao} detalha a implementação computacional da ferramenta \fundaai{}. A Seção~\ref{sec:resultados} reporta os resultados do protocolo experimental. A Seção~\ref{sec:discussao} discute esses resultados de forma integrada com as referências da literatura. Por fim, a Seção~\ref{sec:conclusoes} sintetiza as conclusões e aponta os trabalhos futuros, com destaque para a frente de empacotamento.
